% Dodatak B: Upotreba velikih jezičkih modela u izradi rada
\chapter{Upotreba velikih jezičkih modela u izradi rada}
\label{app:llm}
\chaptermark{Upotreba jezičkih modela u izradi rada}

Tokom izrade ovog master rada korišćeni su veliki jezički modeli kao pomoćni alat u pisanju rada
i razvoju softverskog rešenja predstavljenog u radu.

\paragraph{Upotreba pri pisanju rada.}
Veliki jezički modeli korišćeni su za jezičko i stilsko unapređivanje pojedinih delova teksta,
preformulisanje rečenica radi veće jasnoće i čitljivosti, kao i za pomoć pri organizaciji i strukturiranju sadržaja.
Modeli nisu korišćeni kao zamena za stručnu i naučnu literaturu.
Autor je proverio korišćene izvore i preuzeo odgovornost za tačnost tvrdnji, interpretacija rezultata i zaključaka iznetih u radu.

\paragraph{Upotreba pri razvoju projekta.}
Veliki jezički modeli povremeno su korišćeni kao pomoćni alat tokom razvoja softverskog sistema,
prvenstveno za konsultacije u vezi sa pojedinim implementacionim pitanjima,
razjašnjavanje mogućih uzroka grešaka i predloge za manje izmene ili unapređenja postojećeg koda.
Svi predlozi dobijeni na ovaj način razmotreni su, prilagođeni konkretnim potrebama sistema i provereni od strane autora.
Dizajn sistema, izbor tehnologija, implementacija ključnih funkcionalnosti i testiranje razvijenog rešenja bili su odgovornost autora.

\paragraph{Upotreba u eksperimentalnom delu.}
Pored upotrebe kao pomoćnog alata tokom izrade rada i projekta,
veliki jezički modeli predstavljaju i predmet istraživanja ovog rada i korišćeni su u eksperimentalnoj evaluaciji sistema EduGrader.
Korišćeni modeli, način njihove primene, zadavanje upita i postupak evaluacije detaljno su opisani u odgovarajućim poglavljima rada
(poglavlja~\ref{sec:pipeline}, \ref{sec:metodologija} i~\ref{sec:rezultati}, dodatak~\ref{app:promptovi}).

\paragraph{Korišćeni alati.}
U procesu izrade rada i razvoja projekta korišćeni su:
ChatGPT (OpenAI) za konsultacije i razmatranje ideja tokom razvoja;
\emph{GitHub Copilot}, u razvojnom okruženju \emph{VS Code}, za dopunjavanje linija koda pri kucanju;
Claude (Anthropic), kroz alat \emph{Claude Code}, za jezičko i stilsko uređivanje teksta rada i završno sređivanje korisničkog interfejsa aplikacije.

\paragraph{Odgovornost autora.}
Veliki jezički modeli korišćeni su isključivo kao pomoćni alati.
Autor je samostalno donosio odluke o sadržaju rada i dizajnu i implementaciji sistema,
proveravao i prilagođavao generisane predloge
i izvršio konačnu proveru teksta, programskog koda, eksperimentalnih rezultata i zaključaka rada.
